% =============================
% EINLEITUNG
% =============================

\refstepcounter{chapter}
\chapter*{\thechapter. Einleitung}
\addcontentsline{toc}{chapter}{\thechapter.\space Einleitung}

Die rasante Entwicklung der Quantencomputertechnologie stellt die Informationssicherheit vor grundlegende Herausforderungen. 
Während klassische Computer an physikalische Grenzen stossen, versprechen Quantencomputer durch die Nutzung quantenmechanischer 
Phänomene wie Superposition und Verschränkung für bestimmte mathematische Probleme eine signifikant gesteigerte Rechenleistung. 
Diese technologische Revolution birgt jedoch ein erhebliches Risiko für die Sicherheit heutiger Verschlüsselungsverfahren, die auf mathematischen Problemen basieren.
Besonders betroffen sind asymmetrische Verschlüsselungsverfahren wie RSA, sie könnten durch ausreichend leistungsfähige Quantencomputer vollständig entschlüsselt werden \cite{ct2025verschluesselung}.
\\
Die vorliegende Berufsmaturitätsarbeit untersucht die konkrete Bedrohungslage, die von Quantencomputern für aktuelle Verschlüsselungsstandards ausgeht. 
Im Zentrum steht die Fragestellung: \textit{Welche Bedrohung stellen Quantencomputer für die Sicherheit aktueller Verschlüsselungsverfahren dar, und inwiefern wären dadurch Daten gefährdet, wenn morgen ein Quantencomputer marktreif wäre?}
\\
Zur Beantwortung dieser Frage werden zunächst die physikalischen Grundlagen der Quantencomputertechnologie erläutert, um die Funktionsweise von Qubits und deren Überlegenheit gegenüber klassischen Bits zu verstehen. 
Anschliessend erfolgt eine detaillierte Analyse moderner Verschlüsselungsverfahren, wobei der Fokus auf der RSA-Verschlüsselung und deren Verwundbarkeit gegenüber quantenbasierten Angriffsalgorithmen wie dem Shor-Algorithmus liegt. 
Als messbare Kriterien dienen die erforderliche Anzahl logischer Quantenbits (Qubits) zur Entschlüsselung gängiger Schlüssellängen sowie die Zeitspanne bis zur voraussichtlichen Marktreife derartiger Systeme.
\\
Methodisch stützt sich diese Arbeit auf Fachliteratur und Experteninterviews. 
Ergänzend wurde ein qualitatives Interview mit Dipl.-Ing. Christian Weber MBA, einem IT-Sicherheitsexperten an der ZHAW, geführt. 
Diese Primärquelle ermöglicht eine praxisnahe Einschätzung der aktuellen Risiken.
\\
Die wissenschaftliche Literatur zeigt, dass das Thema in der Fachwelt intensiv diskutiert wird, insbesondere die Bedrohung für bestehende Verschlüsselungsverfahren. 
Die Aktualität der Bedrohung wird durch die konkreten Roadmaps führender Unternehmen wie IBM unterstrichen, die die Entwicklung von Quantencomputern mit über 2000 logischen Qubits für die 2030er Jahre prognostizieren \cite{ct2025verschluesselung}.
